\editionPage[\pdfbookmark[0]{Corrigés des défis}{bm-defi-cor-1-1-parent}]{defi-cor-1-1}{D01 · Corrigé}{Décimales finies et périodes}{\hyperref[defi-ex-1]{Énoncé p.\,\pageref*{defi-ex-1}} · Démarche et contrôles}
\begin{blocEdition}{defi-cor-1-1}
\editionEtape{1. Caractériser les fractions décimales}
Si $p/q=N/10^k$, alors $p10^k=Nq$. Comme $p$ et $q$ sont premiers entre eux, tous les facteurs premiers de $q$ doivent diviser $10^k$ : ce sont 2 et 5. Réciproquement, si $q=2^a5^b$, choisir $k\ge a,b$ permet de multiplier numérateur et dénominateur pour obtenir $10^k$. La fraction a donc une écriture décimale finie.
\editionEtape{2. Calculer une période}
Pour $1/7$, les chiffres produits sont 1, 4, 2, 8, 5, 7 et les restes successifs 3, 2, 6, 4, 5, 1. Le reste initial 1 revient : $\res{1/7=0{,}\overline{142857}}$.
\editionEtape{3. Justifier la répétition}
S’il apparaît un reste nul, la division s’arrête. Sinon, les restes non nuls sont en nombre fini : l’un finit par revenir. À partir d’un même reste, la division produit les mêmes chiffres et les mêmes restes. Le développement devient périodique.
\editionRetenir{Le nombre fini de restes explique la périodicité. Observer six chiffres ne suffit pas sans le retour d’un reste.}
\end{blocEdition}
\editionPage{defi-cor-2-1}{D02 · Corrigé}{Deux identités à construire}{\hyperref[defi-ex-2]{Énoncé p.\,\pageref*{defi-ex-2}} · Démarche et contrôles}
\begin{blocEdition}{defi-cor-2-1}
\editionEtape{1. Compter tous les produits}
La multiplication des deux sommes donne $a^2+b^2+c^2+ab+ba+ac+ca+bc+cb$, soit
\[\res{(a+b+c)^2=a^2+b^2+c^2+2ab+2ac+2bc}.\]
\editionEtape{2. Construire le cube}
On multiplie $a^2+2ab+b^2$ par $a+b$. En regroupant :
\[\res{(a+b)^3=a^3+3a^2b+3ab^2+b^3}.\]
\editionEtape{3. Calculer avec un choix adapté}
Pour $103^3$, choisir $a=100$, $b=3$ : $1000000+90000+2700+27=\res{1092727}$. Pour le carré à trois termes, on obtient $100+4+1+40+20+4=\res{169}$.
\editionRetenir{Les coefficients proviennent du nombre d’occurrences de chaque produit. Une formule se reconstruit si l’on maîtrise la distributivité.}
\end{blocEdition}
\editionPage{defi-cor-3-1}{D03 · Corrigé}{Moyennes et rectangle optimal}{\hyperref[defi-ex-3]{Énoncé p.\,\pageref*{defi-ex-3}} · Démarche et contrôles}
\begin{blocEdition}{defi-cor-3-1}
\editionEtape{1. Établir l’inégalité}
Le carré donne $a+b-2\sqrt{ab}\ge0$, donc $\res{\sqrt{ab}\le(a+b)/2}$. L’égalité exige $\sqrt a=\sqrt b$, donc $a=b$.
\editionEtape{2. Maximiser à somme fixée}
Si $a+b=18$, alors $\sqrt{ab}\le9$, donc $ab\le81$. L’égalité est atteinte pour $\res{a=b=9}$ : l’aire maximale vaut $\res{81}$.
\editionEtape{3. Minimiser à produit fixé}
Si $ab=36$, alors $(a+b)/2\ge6$, donc le périmètre $2(a+b)\ge24$. La borne est atteinte pour $\res{a=b=6}$. Le périmètre minimal vaut $\res{24}$.
\editionRetenir{Une inégalité n’établit un optimum qu’après vérification que son cas d’égalité est réalisable sous la contrainte.}\label{fig-defi-cor-3}
\figureGeo[corrige]{defi-moyennes}

\end{blocEdition}
\editionPage{defi-cor-4-1}{D04 · Corrigé}{Barycentres de deux et trois points}{\hyperref[defi-ex-4]{Énoncé p.\,\pageref*{defi-ex-4}} · Démarche et contrôles}
\begin{blocEdition}{defi-cor-4-1}
\editionEtape{1. Résoudre l’égalité vectorielle}
En coordonnées, $\alpha(A-G)+\beta(B-G)=0$ donne $(\alpha+\beta)G=\alpha A+\beta B$. Ainsi
\[G\left(\frac{\alpha x_A+\beta x_B}{\alpha+\beta};\frac{\alpha y_A+\beta y_B}{\alpha+\beta}\right).\]
Pour les poids 2 et 1, $\res{G(2;0)}$. Pour 2 et $-1$, $\res{G(-6;0)}$ ; le point sort bien du segment.
\editionEtape{2. Étendre à trois points}
Le même regroupement donne $G=(\alpha A+\beta B+\gamma C)/(\alpha+\beta+\gamma)$. Avec les poids 1, 1 et 2, on obtient $\res{G(3/2;3)}$. On contrôle que les trois vecteurs pondérés depuis ce point ont une somme nulle.
\editionRetenir{La somme des poids doit être non nulle pour déterminer ce point par la formule. Le milieu correspond aux deux poids égaux.}\label{fig-defi-cor-4}
\figureGeo[corrige]{defi-barycentres}

\end{blocEdition}
\editionPage{defi-cor-5-1}{D05 · Corrigé}{Pourquoi les médianes se coupent}{\hyperref[defi-ex-5]{Énoncé p.\,\pageref*{defi-ex-5}} · Démarche et contrôles}
\begin{blocEdition}{defi-cor-5-1}
\editionEtape{1. Comparer les deux déplacements}
On a $G-A=(B+C-2A)/3$ et $I-A=(B+C-2A)/2$. Donc $\res{\vect{AG}=\tfrac23\vect{AI}}$. Cela place $G$ sur la médiane issue de $A$, aux deux tiers depuis le sommet.
\editionEtape{2. Passer aux trois médianes}
Le même calcul en échangeant les lettres place $G$ sur les médianes issues de $B$ et de $C$. Les médianes sont donc concourantes. Deux médianes d’un triangle non aplati sont distinctes, leur intersection est unique.
\editionEtape{3. Contrôler l’exemple}
Pour le triangle proposé, $\res{G(2;2)}$. Les milieux sont $(3;3)$, $(0;3)$ et $(3;0)$. Les déplacements depuis chaque sommet confirment le coefficient $2/3$.
\editionRetenir{L’égalité vectorielle donne à la fois l’appartenance aux médianes et la position précise du point de concours.}\label{fig-defi-cor-5}
\figureGeo[corrige]{defi-medianes}

\end{blocEdition}
\editionPage{defi-cor-6-1}{D06 · Corrigé}{Concours des hauteurs}{\hyperref[defi-ex-6]{Énoncé p.\,\pageref*{defi-ex-6}} · Démarche et contrôles}
\begin{blocEdition}{defi-cor-6-1}
\editionEtape{1. Situer les milieux}
$B'+C'=(A+C-B)+(A+B-C)=2A$. Donc $A$ est le milieu de $[B'C']$. De même, $B$ et $C$ sont les milieux des deux autres côtés.
\editionEtape{2. Relier les directions}
$\vect{B'C'}=C'-B'=2(B-C)=-2\vect{BC}$. Les côtés correspondants sont parallèles. La droite passant par $A$ et perpendiculaire à $(BC)$ est donc aussi perpendiculaire à $(B'C')$ en son milieu : c’est sa médiatrice.
\editionEtape{3. Utiliser un concours déjà établi}
Les trois hauteurs de $ABC$ sont les trois médiatrices du triangle $A'B'C'$. Celles-ci sont concourantes par D08 ; les hauteurs le sont donc également. Leur point de concours est appelé orthocentre.
\editionRetenir{Transformer une configuration en une autre peut ramener une preuve à un théorème déjà connu. Le recours à D08 est annoncé, pour éviter une preuve circulaire.}\label{fig-defi-cor-6}
\figureGeo[corrige]{defi-hauteurs}

\end{blocEdition}
\editionPage{defi-cor-7-1}{D07 · Corrigé}{Une aire à partir de deux côtés et d’un angle}{\hyperref[defi-ex-7]{Énoncé p.\,\pageref*{defi-ex-7}} · Démarche et contrôles}
\begin{blocEdition}{defi-cor-7-1}
\editionEtape{1. Identifier la hauteur}
Avec la base $CB=a$, la hauteur issue de $A$ vaut $h=b\sin C$ lorsque l’angle est aigu. Donc $\mathcal A=ah/2=\res{ab\sin C/2}$.
\editionEtape{2. Étendre au cas obtus}
La hauteur tombe alors sur le prolongement de la base. Dans le triangle rectangle auxiliaire, elle vaut $b\sin(180^\circ-C)$. La définition annoncée du sinus obtus donne la même formule. Le cas droit donne directement $ab/2$.
\editionEtape{3. Appliquer sans surinterpréter}
Pour 30 et 150 degrés, le sinus vaut $1/2$. Les deux aires valent $\res{10}$. Les troisièmes côtés ne sont pas les mêmes : l’angle et la position du sommet changent. Une aire identique ne détermine pas un triangle unique.
\editionRetenir{L’extension du sinus aux angles obtus est explicitée ici ; elle n’est pas supposée acquise dans le cours de base.}\label{fig-defi-cor-7}
\figureGeo[corrige]{defi-aire-angle}

\end{blocEdition}
\editionPage{defi-cor-8-1}{D08 · Corrigé}{Concours des médiatrices et cercle circonscrit}{\hyperref[defi-ex-8]{Énoncé p.\,\pageref*{defi-ex-8}} · Démarche et contrôles}
\begin{blocEdition}{defi-cor-8-1}
\editionEtape{1. Assurer l’intersection}
Les droites $(AB)$ et $(AC)$ sont sécantes et distinctes puisque le triangle n’est pas aplati. Leurs perpendiculaires ne sont donc pas parallèles. Les deux médiatrices se coupent en un unique point $O$.
\editionEtape{2. Établir le troisième rattachement}
L’appartenance aux deux médiatrices donne $OA=OB$ et $OA=OC$, donc $OB=OC$. Par la caractérisation de la médiatrice, $O$ est aussi sur celle de $[BC]$.
\editionEtape{3. Construire et vérifier le cercle}
Le cercle de centre $O$ et de rayon $OA>0$ passe par $A,B,C$. Tout autre centre possible serait sur les deux premières médiatrices, donc égal à $O$. Dans l’exemple, $x_O=3$. L’égalité $OA^2=OC^2$ donne $9+y^2=1+(y-4)^2$, d’où $\res{O(3;1)}$ et $\res{R=\sqrt{10}}$.
\editionRetenir{Le centre d’un cercle circonscrit n’est pas toujours à l’intérieur du triangle. La preuve repose sur l’équidistance, indépendamment du dessin.}\label{fig-defi-cor-8}
\figureGeo[corrige]{defi-circonscrit}

\end{blocEdition}
\editionPage{defi-cor-9-1}{D09 · Corrigé}{Un lieu parabolique}{\hyperref[defi-ex-9]{Énoncé p.\,\pageref*{defi-ex-9}} · Démarche et contrôles}
\begin{blocEdition}{defi-cor-9-1}
\editionEtape{1. Traduire et simplifier}
$MF^2=x^2+(y-2)^2$ et la distance à l’axe a pour carré $y^2$. L’égalité donne $x^2+y^2-4y+4=y^2$, donc $\res{y=1+x^2/4}$.
\editionEtape{2. Contrôler la réciproque}
Pour un point vérifiant cette équation, le calcul inverse donne $MF^2=y^2$. Les deux distances étant positives, $MF=|y|$. Le point appartient bien au lieu. L’équation obtenue décrit ainsi exactement l’ensemble, sans point parasite.
\editionEtape{3. Exploiter le modèle}
Le carré est positif, donc $y\ge1$, avec égalité pour $x=0$ : le point le plus bas est $\res{(0;1)}$. Sur $y=5$, $x^2=16$ ; les points sont $\res{(-4;5)\text{ et }(4;5)}$.
\editionRetenir{Un ensemble défini par une distance peut devenir une courbe d’équation simple. Il faut prouver les deux sens de l’équivalence.}\label{fig-defi-cor-9}
\figureGeo[corrige]{defi-parabole}

\end{blocEdition}
\editionPage{defi-cor-10-1}{D10 · Corrigé}{Des inégalités décrivent une région}{\hyperref[defi-ex-10]{Énoncé p.\,\pageref*{defi-ex-10}} · Démarche et contrôles}
\begin{blocEdition}{defi-cor-10-1}
\editionEtape{1. Construire le domaine}
Les sommets sont $\res{(0;0),(8;0),(4;4),(0;6)}$. Chaque inégalité sélectionne un côté d’une droite ; leur intersection donne le polygone, bords inclus.
\editionEtape{2. Prouver une borne sur toute la région}
On écrit $3x+4y=2(x+y)+(x+2y)$. Avec les deux contraintes, $R\le2\times8+12=\res{28}$. Ce calcul est valable pour tous les points admissibles, pas seulement pour les sommets.
\editionEtape{3. Vérifier l’atteinte}
Pour avoir égalité, il faut $x+y=8$ et $x+2y=12$, d’où $\res{(x;y)=(4;4)}$. Ce point respecte les contraintes de positivité. Le maximum est donc 28 et il est atteint uniquement là.
\editionRetenir{Une combinaison d’inégalités peut prouver un optimum sur une région entière, sans théorème avancé d’optimisation.}\label{fig-defi-cor-10}
\figureGeo[corrige]{defi-region}

\end{blocEdition}
\editionPage{defi-cor-11-1}{D11 · Corrigé}{Carré et racine, deux courbes liées}{\hyperref[defi-ex-11]{Énoncé p.\,\pageref*{defi-ex-11}} · Démarche et contrôles}
\begin{blocEdition}{defi-cor-11-1}
\editionEtape{1. Échanger les coordonnées}
Pour $a\ge0$, $\sqrt{a^2}=a$. Le point $(a^2;a)$ appartient donc à la courbe racine. Échanger les coordonnées décrit la symétrie par rapport à la droite $y=x$. Il faut restreindre la fonction carré à $[0;+\infty[$ : sa branche négative ne devient pas la courbe de la racine positive.
\editionEtape{2. Trouver les intersections}
$\sqrt x=x^2$ impose $x\ge0$. En élevant au carré, $x=x^4$, donc $x(x^3-1)=0$. La croissance du cube donne $\res{x=0\text{ ou }x=1}$. Les points sont $(0;0)$ et $(1;1)$.
\editionEtape{3. Comparer les trois fonctions}
Pour $0<x<1$, $\res{x^2<x<\sqrt x}$. Pour $x>1$, $\res{\sqrt x<x<x^2}$. Les inégalités sont strictes hors des deux points d’égalité.
\editionRetenir{La symétrie n’est valable qu’après choix de la branche où la fonction carré est inversible.}\label{fig-defi-cor-11}
\figureGeo[corrige]{defi-racine-carre}

\end{blocEdition}
\editionPage{defi-cor-12-1}{D12 · Corrigé}{Un paramètre et deux extrémums}{\hyperref[defi-ex-12]{Énoncé p.\,\pageref*{defi-ex-12}} · Démarche et contrôles}
\begin{blocEdition}{defi-cor-12-1}
\editionEtape{1. Classer les positions du centre}
Si $a<0$, le plus proche est 0 et le minimum vaut $a^2$. Si $0\le a\le4$, il vaut 0, atteint en $x=a$. Si $a>4$, il vaut $(4-a)^2$, atteint en 4.
\editionEtape{2. Comparer les deux extrémités}
Le maximum est $\max(a^2,(4-a)^2)$. Leur différence vaut $a^2-(4-a)^2=8a-16$. Le premier est plus grand pour $a>2$, le second pour $a<2$, et ils sont égaux pour $a=2$.
\editionEtape{3. Optimiser la pire valeur}
Si $a\le2$, alors $4-a\ge2$, donc le maximum est au moins 4. Si $a\ge2$, alors $a\ge2$, même conclusion. La valeur 4 est atteinte pour $\res{a=2}$ et seulement pour ce choix ; le maximum minimal vaut $\res4$.
\editionRetenir{Un paramètre oblige à distinguer les régimes. Un optimum portant sur une pire valeur se traite encore avec des comparaisons élémentaires.}
\end{blocEdition}
\editionPage{defi-cor-13-1}{D13 · Corrigé}{Optimiser une dispersion}{\hyperref[defi-ex-13]{Énoncé p.\,\pageref*{defi-ex-13}} · Démarche et contrôles}
\begin{blocEdition}{defi-cor-13-1}
\editionEtape{1. Écrire les indicateurs}
La moyenne est $\bar x=(t+6)/3$ et la médiane est $t$. La moyenne des carrés vaut $(t^2+36)/3$.
\editionEtape{2. Transformer la dispersion}
On calcule $\sigma^2=(t^2+36)/3-(t+6)^2/9=(2t^2-12t+72)/9$, soit
\[\res{\sigma^2=\frac29(t-3)^2+6}.\]
La racine étant croissante, minimiser $\sigma^2$ revient à minimiser $\sigma$.
\editionEtape{3. Conclure avec les conditions}
Le minimum est atteint pour $\res{t=3}$, autorisé par le domaine ; l’écart type vaut alors $\res{\sqrt6}$. L’équation $\bar x=3$ donne aussi $(t+6)/3=3$, donc le même choix.
\editionRetenir{La moyenne, la médiane et l’écart type décrivent des propriétés distinctes. Une forme carrée peut rendre une optimisation statistique accessible.}
\end{blocEdition}
\editionPage{defi-cor-14-1}{D14 · Corrigé}{Des marges qui ne fixent pas l’intersection}{\hyperref[defi-ex-14]{Énoncé p.\,\pageref*{defi-ex-14}} · Démarche et contrôles}
\begin{blocEdition}{defi-cor-14-1}
\editionEtape{1. Reconstituer le tableau}
Les quatre cases sont $k$ pour les deux, $60-k$ pour $A$ seul, $50-k$ pour $B$ seul et $k-10$ pour aucun des deux. Leur somme vaut bien 100.
\editionEtape{2. Imposer la faisabilité}
La positivité des cases exige $k\ge10$ et $k\le50$. Toute valeur entière de $\res{10\text{ à }50}$ produit un tableau valide ; ces bornes ne sont pas seulement nécessaires, elles sont réalisables.
\editionEtape{3. Encadrer le conditionnement}
La fréquence cherchée est $k/60$. Elle peut prendre les valeurs discrètes correspondantes entre $\res{1/6\text{ et }5/6}$. Les marges ne permettent donc pas de connaître une unique fréquence conditionnelle.
\editionRetenir{Un tableau n’est pas déterminé par ses seules marges. Vérifier les quatre cases permet aussi de détecter une intersection impossible.}
\end{blocEdition}
\editionPage{defi-cor-15-1}{D15 · Corrigé}{Défauts rares et conditionnement inversé}{\hyperref[defi-ex-15]{Énoncé p.\,\pageref*{defi-ex-15}} · Démarche et contrôles}
\begin{blocEdition}{defi-cor-15-1}
\editionEtape{1. Organiser les deux origines du signal}
Les vrais positifs sont $0{,}9pN$ ; les faux positifs $0{,}05(1-p)N$. Leur somme donne tous les signaux positifs. Pour un modèle de proportions, ces expressions n’imposent pas que chaque taille arbitraire $N$ donne des effectifs entiers.
\editionEtape{2. Restreindre la référence aux positifs}
Le rapport vrais positifs sur tous les positifs simplifie le facteur $N$ et donne la fraction annoncée.
\editionEtape{3. Établir le seuil}
La fraction est au moins $1/2$ si $0{,}9p\ge0{,}05(1-p)$, soit $0{,}95p\ge0{,}05$. Donc $\res{p\ge1/19\approx5{,}26\,\%}$. Pour $p=1\,\%$, la probabilité inversée vaut $2/13\approx15{,}4\,\%$ ; pour $p=10\,\%$, elle vaut $2/3\approx66{,}7\,\%$.
\editionRetenir{Les performances du contrôle ne suffisent pas à connaître la probabilité inversée : la rareté initiale des défauts intervient aussi.}\label{fig-defi-cor-15}
\figureGeo[corrige]{defi-test}

\end{blocEdition}
\editionPage{defi-cor-16-1}{D16 · Corrigé}{Trois dés, deux sommes}{\hyperref[defi-ex-16]{Énoncé p.\,\pageref*{defi-ex-16}} · Démarche et contrôles}
\begin{blocEdition}{defi-cor-16-1}
\editionEtape{1. Distinguer les types et les issues}
Pour 9 : $(1,2,6),(1,3,5),(1,4,4),(2,2,5),(2,3,4),(3,3,3)$, avec multiplicités 6, 6, 3, 3, 6, 1 : total 25. Pour 10 : $(1,3,6),(1,4,5),(2,2,6),(2,3,5),(2,4,4),(3,3,4)$, avec 6, 6, 3, 6, 3, 3 : total 27.
\editionEtape{2. Calculer les probabilités}
Les $6^3=216$ triplets ordonnés sont équiprobables. Donc $\res{P(S=9)=25/216}$ et $\res{P(S=10)=27/216=1/8}$. La différence est $\res{1/108}$ en faveur de 10. Les six types non ordonnés n’avaient pas tous la même probabilité.
\editionRetenir{Il faut compter des issues équiprobables, et non des descriptions qui regroupent des nombres différents d’issues.}
\end{blocEdition}
\editionPage{defi-cor-16-2}{D16 · Corrigé (suite)}{Trois dés, deux sommes}{\hyperref[defi-ex-16]{Énoncé p.\,\pageref*{defi-ex-16}} · Démarche et contrôles}
\begin{blocEdition}{defi-cor-16-2}
\editionEtape{3. Vérifier par une énumération exhaustive}
\begin{python}
def effectifs():
    n9, n10 = 0, 0
    for a in range(1, 7):
        for b in range(1, 7):
            for c in range(1, 7):
                s = a + b + c
                if s == 9:
                    n9 += 1
                if s == 10:
                    n10 += 1
    return n9, n10
\end{python}
La fonction renvoie \code{(25, 27)}. Chaque triplet ordonné est visité exactement une fois, car chaque boucle parcourt toutes les valeurs de son dé. Il s’agit d’une énumération exhaustive, pas d’une simulation aléatoire : le dénombrement est donc exact.
\editionRetenir{L’informatique peut servir à prouver un énoncé fini lorsqu’on justifie que tous les cas sont examinés exactement une fois.}
\end{blocEdition}
